% INDEPENDENT INSTRUCTIONAL MATERIAL --- NOT WEN-WEI LI'S SOURCE TEXT.
% stable-unit-id: o014.mastery.derived-spectral.002
% provenance: OpenAI Codex gpt-5.6-sol, Ultra
% license: CC BY 4.0; compatible with the source translation license.

\section*{Mastery Bridge 002: Derived Functors and Spectral Sequences}
\addcontentsline{toc}{section}{Mastery Bridge 002: Derived Functors and Spectral Sequences}
\hypertarget{o014.mastery.derived-spectral.002}{ }
\label{mastery:derived-spectral-002}

\noindent\textbf{Status and attribution.}
This unit is independent instructional material prepared to accompany this
book. It is not part of Wen-Wei Li's source text and must not be attributed to
him. Production provenance:
\texttt{OpenAI Codex gpt-5.6-sol, Ultra}. This material is available under the
Creative Commons Attribution 4.0 International license
(CC BY 4.0; \url{https://creativecommons.org/licenses/by/4.0/}), and is
therefore compatible with the license of the source text.

\subsection*{Prerequisite map}

The shortest prerequisite chain is as follows.
\begin{compactitem}
	\item Abelian categories, complexes, and cohomology:
	\S\ref{sec:Abel-cat-def}, \S\ref{sec:cohomology}, and
	Definition~\ref{def:cohomology}.
	\item Injective/projective objects and the existence of enough such
	objects:
	\S\ref{sec:inj-proj}, Definition~\ref{def:injective-projective-obj}, and
	Definition~\ref{def:enough-injectives}.
	\item Resolutions and uniqueness up to homotopy:
	\S\ref{sec:resolutions}, Definition~\ref{def:resolutions}, and
	Theorem~\ref{prop:resolution-homotopy}.
	\item Classical derived functors, long exact sequences, and dimension
	shifting: \S\ref{sec:derived-primer},
	Theorem~\ref{prop:long-exact-sequence-primer}, and
	Proposition~\ref{prop:dimension-shifting}.
	\item Spectral sequences, edge maps, and composition of functors:
	\S\ref{sec:spectral-sequence}, Remark~\ref{rem:edge-computing}, and
	Theorem~\ref{prop:Grothendieck-ss}.
	\item Cohomology of cyclic groups and multiplicative structure:
	Definition~\ref{def:group-coh-ho},
	Proposition~\ref{prop:group-cohomology-cyclic},
	\S\ref{sec:grading-filtration}, and
	\S\ref{sec:multiplicative-SS}.
\end{compactitem}
The conceptual progression is
\[
	\text{replace by a resolution}
	\Longrightarrow \text{take cohomology}
	\Longrightarrow \text{organize the total filtration}
	\Longrightarrow \text{read the pages and edges}.
\]

\subsection*{Working method: three ledgers}

To keep a computation organized, record its information in three separate
``ledgers'' that must not be conflated.
\begin{enumerate}
	\item The \emph{resolution ledger} records which object is replaced by a
	resolution, the direction of that resolution, and why its terms are
	injective, projective, or acyclic for the functor to be applied.
	\item The \emph{degree ledger} records cohomological indices, total degree
	$n=p+q$, the direction of the differential
	$d_r:E_r^{p,q}\to E_r^{p+r,q-r+1}$, and every vanishing term.
	\item The \emph{filtration ledger} records the abutment object $H^n$, the
	filtration $\mathrm{F}^pH^n$, the identification
	$\gr^pH^n\simeq E_\infty^{p,n-p}$, and multiplicative compatibility.
\end{enumerate}
The most common mistake is to identify $E_2^{p,q}$ directly with a subobject
of the abutment. What is canonically available is the graded term
$E_\infty^{p,q}$; returning to the abutment still requires solving the
filtration extension problem.

As a worked example, take a first-quadrant cohomological spectral sequence
\[
	E_2^{p,q}\Rightarrow H^{p+q}.
\]
In total degree $1$, no differential touches $E_2^{1,0}$, while the only
differential that can leave $E_2^{0,1}$ is
$d_2:E_2^{0,1}\to E_2^{2,0}$. Therefore
\[
	E_\infty^{1,0}=E_2^{1,0},\qquad
	E_\infty^{0,1}=\Ker(d_2:E_2^{0,1}\to E_2^{2,0}).
\]
The filtration of $H^1$ has a subobject
$\mathrm{F}^1H^1\simeq E_\infty^{1,0}$ and quotient
$H^1/\mathrm{F}^1H^1\simeq E_\infty^{0,1}$. In total degree $2$,
\[
	E_\infty^{2,0}\simeq
	E_2^{2,0}/\Image(d_2:E_2^{0,1}\to E_2^{2,0})
	\hookrightarrow H^2.
\]
These three observations immediately produce the five-term sequence
\[
	0\to E_2^{1,0}\to H^1\to E_2^{0,1}
	\xrightarrow{d_2}E_2^{2,0}\to H^2.
\]
The method is always the same: mark the vanishing region, determine which
differentials remain possible, compute $E_\infty$, and only then use the
abutment filtration.

\subsection*{Exercises and complete solutions}

% stable-exercise-id: o014.mastery.derived-spectral.002.ex01
\subsubsection*{Exercise 1 --- Comparing resolutions and independence of choices}
\hypertarget{o014.mastery.derived-spectral.002.ex01}{ }
\label{mastery:derived-spectral-002-ex01}

Let $\mathcal{A}$ be an abelian category with enough injectives, let
$F:\mathcal{A}\to\mathcal{B}$ be an additive functor, and let
\[
	0\to X\to I^0\to I^1\to\cdots,
	\qquad
	0\to X\to J^0\to J^1\to\cdots
\]
be two injective resolutions of $X$.
\begin{enumerate}[(i)]
	\item Construct cochain maps $u:I^\bullet\to J^\bullet$ and
	$v:J^\bullet\to I^\bullet$ lifting $\identity_X$.
	\item Prove by a cochain homotopy that $vu\simeq\identity_I$ and
	$uv\simeq\identity_J$.
	\item Conclude that $\Hm^n(FI^\bullet)$ and
	$\Hm^n(FJ^\bullet)$ are canonically isomorphic for every $n$.
\end{enumerate}

% stable-solution-id: o014.mastery.derived-spectral.002.sol01
\paragraph{Solution.}
\hypertarget{o014.mastery.derived-spectral.002.sol01}{ }
Because $X\hookrightarrow I^0$ is a monomorphism and $J^0$ is injective,
the morphism $X\to J^0$ extends to $u^0:I^0\to J^0$. Suppose that
$u^0,\ldots,u^n$ have been constructed and satisfy the cochain-map
condition. The morphism $d_J^nu^n:I^n\to J^{n+1}$ vanishes on
$\Ker(d_I^n)=\Image(d_I^{n-1})$, and hence factors through
$\Image(d_I^n)\hookrightarrow I^{n+1}$. Since $J^{n+1}$ is injective,
this factorization extends to $u^{n+1}:I^{n+1}\to J^{n+1}$ with
\[
	u^{n+1}d_I^n=d_J^nu^n.
\]
Induction produces $u$. Interchanging $I$ and $J$ produces $v$.

We also need the fact that any two lifts $r,s:I^\bullet\to J^\bullet$ of
$\identity_X$ are homotopic. In degree zero, $r^0-s^0$ vanishes on
$X=\Ker(d_I^0)$, and hence factors through
$\Image(d_I^0)\hookrightarrow I^1$. Injectivity of $J^0$ gives
$h^1:I^1\to J^0$ with
$r^0-s^0=h^1d_I^0$. Inductively, once $h^1,\ldots,h^n$ have been
constructed, the morphism
\[
	r^n-s^n-d_J^{n-1}h^n:I^n\longrightarrow J^n
\]
vanishes on $\Ker(d_I^n)$. It factors through
$\Image(d_I^n)\hookrightarrow I^{n+1}$ and, because $J^n$ is injective,
extends to $h^{n+1}:I^{n+1}\to J^n$. Thus
\[
	r-s=d_Jh+hd_I,
\]
that is, $r\simeq s$.

The composites $vu$ and $\identity_I$ both lift $\identity_X$, so they are
homotopic; similarly, $uv\simeq\identity_J$. Since $F$ is additive, it
carries the homotopy equations to homotopy equations. Therefore
$\Hm^n(Fu)$ and $\Hm^n(Fv)$ are mutually inverse. If other comparison maps
are chosen, they are homotopic to the first ones and induce the same maps on
cohomology. The isomorphism is consequently canonical. This is the concrete
content of independence of resolutions in
Theorem~\ref{prop:resolution-homotopy}.

% stable-exercise-id: o014.mastery.derived-spectral.002.ex02
\subsubsection*{Exercise 2 --- Computing $\Ext$ from a short resolution}
\hypertarget{o014.mastery.derived-spectral.002.ex02}{ }
\label{mastery:derived-spectral-002-ex02}

Let $m\geq1$ and let $M$ be an abelian group. Use the projective resolution
\[
	0\longrightarrow\Z\xrightarrow{\times m}\Z
	\longrightarrow\Z/m\Z\longrightarrow0
\]
to compute $\Ext^n_{\Z}(\Z/m\Z,M)$ for all $n\geq0$. Explain why the
answer is independent of the chosen projective resolution.

% stable-solution-id: o014.mastery.derived-spectral.002.sol02
\paragraph{Solution.}
\hypertarget{o014.mastery.derived-spectral.002.sol02}{ }
Apply $\Hom_{\Z}(-,M)$ to the resolution. Since
$\Hom_{\Z}(\Z,M)\simeq M$, this gives the cochain complex
\[
	0\longrightarrow
	\underbracket{M}_{\text{degree }0}
	\xrightarrow{\times m}
	\underbracket{M}_{\text{degree }1}
	\longrightarrow0.
\]
Thus
\[
	\Ext^n_{\Z}(\Z/m\Z,M)\simeq
	\begin{cases}
		M[m]:=\{x\in M:mx=0\},&n=0,\\
		M/mM,&n=1,\\
		0,&n\geq2.
	\end{cases}
\]
For $n=0$, this identification is the usual one
$\Hom_{\Z}(\Z/m\Z,M)\simeq M[m]$. Two projective resolutions of
$\Z/m\Z$ are connected by comparison maps unique up to homotopy, the dual
version of Exercise 1. The functor $\Hom_{\Z}(-,M)$ carries these homotopies
to homotopies, so the cohomology is unchanged. Hence the answer above is
independent of the resolution used.

% stable-exercise-id: o014.mastery.derived-spectral.002.ex03
\subsubsection*{Exercise 3 --- The long exact sequence of derived functors}
\hypertarget{o014.mastery.derived-spectral.002.ex03}{ }
\label{mastery:derived-spectral-002-ex03}

Let $F:\mathcal{A}\to\mathcal{B}$ be an additive left-exact functor, and
suppose that $\mathcal{A}$ has enough injectives. For a short exact sequence
\[
	0\longrightarrow X\longrightarrow Y\longrightarrow Z\longrightarrow0,
\]
use compatible injective resolutions to construct the connecting morphisms
\[
	\delta^n:\mathrm{R}^nF(Z)\longrightarrow\mathrm{R}^{n+1}F(X).
\]
Prove that these morphisms are well defined and natural, and that they
produce a long exact sequence.

% stable-solution-id: o014.mastery.derived-spectral.002.sol03
\paragraph{Solution.}
\hypertarget{o014.mastery.derived-spectral.002.sol03}{ }
The injective version of the Horseshoe Lemma gives a diagram of resolutions
whose rows are exact and which is split in every degree:
\[
\begin{tikzcd}
	0 \arrow[r] & X \arrow[r] \arrow[d] & Y \arrow[r] \arrow[d] & Z \arrow[r] \arrow[d] & 0 \\
	0 \arrow[r] & I_X^\bullet \arrow[r] & I_Y^\bullet \arrow[r] & I_Z^\bullet \arrow[r] & 0.
\end{tikzcd}
\]
Because the bottom row is split in every degree, every additive functor
preserves its exactness. Thus
\[
	0\to F(I_X^\bullet)\to F(I_Y^\bullet)
	\to F(I_Z^\bullet)\to0
\]
is a short exact sequence of complexes.

Take a class $[z]\in\Hm^n(FI_Z^\bullet)$ with cocycle representative $z$.
Lift $z$ to $y\in F(I_Y^n)$. Since the image of $d(y)$ in
$F(I_Z^{n+1})$ is $d(z)=0$, there is an $x\in F(I_X^{n+1})$ whose image
is $d(y)$. The cochain map on the left is injective, and $d^2=0$ then shows
that $d(x)=0$. Define
\[
	\delta^n[z]:=[x].
\]
If $y$ is replaced by another lift, the difference comes from
$F(I_X^n)$ and $x$ changes by a coboundary. If $z$ is replaced by
$z+d(w)$, lift $w$ and correct $y$ by the differential of that lift; the
class $[x]$ remains unchanged. Thus $\delta^n$ is well defined.

Exactness is checked by the same chase. For example,
$[y]\in\Hm^n(FI_Y)$ maps to zero in $\Hm^n(FI_Z)$ if and only if, after
correcting $y$ by a coboundary, it comes from a cocycle in $F(I_X^n)$.
Next, $[z]$ lies in the kernel of $\delta^n$ if and only if $x$ is a
coboundary; correcting the lift $y$ by a preimage of that coboundary then
produces a cocycle mapping to $z$. The analogous argument at the next term
completes the exactness check and yields
\[
	\cdots\to\mathrm{R}^nF(X)\to\mathrm{R}^nF(Y)
	\to\mathrm{R}^nF(Z)\xrightarrow{\delta^n}
	\mathrm{R}^{n+1}F(X)\to\cdots.
\]
For a morphism between two short exact sequences, comparison maps of
resolutions carry the choices $z,y,x$ to corresponding choices in the
second diagram. Both orders of operations produce the same class, so
$\delta^n$ is natural. This construction realizes
Theorem~\ref{prop:long-exact-sequence-primer}.

% stable-exercise-id: o014.mastery.derived-spectral.002.ex04
\subsubsection*{Exercise 4 --- Dimension shifting}
\hypertarget{o014.mastery.derived-spectral.002.ex04}{ }
\label{mastery:derived-spectral-002-ex04}

Let $F:\mathcal{A}\to\mathcal{B}$ be a left-exact functor, and suppose
that
\[
	0\longrightarrow X\longrightarrow A\longrightarrow B\longrightarrow0
\]
is exact, with $A$ being $F$-acyclic.
\begin{enumerate}[(i)]
	\item Prove the natural isomorphisms
	\[
		\mathrm{R}^1F(X)\simeq\Coker[FA\to FB],
		\qquad
		\mathrm{R}^nF(X)\simeq\mathrm{R}^{n-1}F(B)\qquad(n\geq2).
	\]
	\item For an injective resolution
	$0\to X\to I^0\to I^1\to\cdots$, define
	$C^0=X$ and $C^{r+1}=\Coker[C^r\to I^r]$. Apply (i) repeatedly to
	obtain a cohomological formula for $\mathrm{R}^nF(X)$.
	\item If $X$ has an injective resolution with $I^r=0$ for $r>d$,
	prove that $\mathrm{R}^nF(X)=0$ for $n>d$.
\end{enumerate}

% stable-solution-id: o014.mastery.derived-spectral.002.sol04
\paragraph{Solution.}
\hypertarget{o014.mastery.derived-spectral.002.sol04}{ }
The long exact sequence from Exercise 3 contains
\[
	FA\longrightarrow FB\longrightarrow\mathrm{R}^1F(X)
	\longrightarrow\underbracket{\mathrm{R}^1F(A)}_{0}
\]
and, for $n\geq2$,
\[
	\underbracket{\mathrm{R}^{n-1}F(A)}_{0}
	\longrightarrow\mathrm{R}^{n-1}F(B)
	\longrightarrow\mathrm{R}^nF(X)
	\longrightarrow\underbracket{\mathrm{R}^nF(A)}_{0}.
\]
Exactness gives both isomorphisms in (i), and their naturality follows from
the naturality of the long exact sequence.

The injective resolution decomposes into short exact sequences
\[
	0\to C^r\to I^r\to C^{r+1}\to0,
	\qquad r\geq0.
\]
Because $I^r$ is injective, it is $F$-acyclic. Therefore, for $n\geq1$,
\[
	\mathrm{R}^nF(X)\simeq\mathrm{R}^{n-1}F(C^1)\simeq\cdots
	\simeq\mathrm{R}^1F(C^{n-1})
	\simeq\Coker[FI^{n-1}\to FC^n].
\]
Left exactness of $F$ identifies $FC^n$ with
$\Ker[FI^n\to FI^{n+1}]$. Hence
\[
	\mathrm{R}^nF(X)\simeq
	\frac{\Ker[FI^n\to FI^{n+1}]}
	{\Image[FI^{n-1}\to FI^n]}
	=\Hm^n(FI^\bullet).
\]
If the resolution ends at $I^d$, the complex $FI^\bullet$ has no terms in
degrees greater than $d$, so its cohomology vanishes in degrees $n>d$. This
also follows by shifting dimension until one reaches the final injective
cosyzygy. It is the computational form of
Proposition~\ref{prop:dimension-shifting}.

% stable-exercise-id: o014.mastery.derived-spectral.002.ex05
\subsubsection*{Exercise 5 --- Setting up the Grothendieck spectral sequence}
\hypertarget{o014.mastery.derived-spectral.002.ex05}{ }
\label{mastery:derived-spectral-002-ex05}

Let
\[
	\mathcal{A}\xrightarrow{F}\mathcal{B}\xrightarrow{G}\mathcal{C}
\]
be additive left-exact functors. Suppose that $\mathcal{A}$ and
$\mathcal{B}$ have enough injectives, and that $F$ maps every injective
object to a $G$-acyclic object. For $X\in\Obj(\mathcal{A})$, prove the
existence of a first-quadrant spectral sequence
\[
	E_2^{p,q}=\mathrm{R}^pG(\mathrm{R}^qF(X))
	\Rightarrow\mathrm{R}^{p+q}(GF)(X).
\]
Explain explicitly the role of the $G$-acyclicity hypothesis.

% stable-solution-id: o014.mastery.derived-spectral.002.sol05
\paragraph{Solution.}
\hypertarget{o014.mastery.derived-spectral.002.sol05}{ }
Choose an injective resolution $0\to X\to I^\bullet$. The complex
$F(I^\bullet)$ lies in nonnegative degrees. Choose a Cartan--Eilenberg
resolution $F(I^\bullet)\to J^{\bullet,\bullet}$ in $\mathcal{B}$ and
apply $G$ to obtain a first-quadrant double complex $GJ^{p,q}$. The two
filtrations of the total complex $\tot(GJ)$ give two spectral sequences
converging to the same abutment.

For the first filtration, take cohomology in the direction that resolves
each $F(I^p)$. This gives
\[
	\Hm^q(GJ^{p,\bullet})
	\simeq\mathrm{R}^qG(F(I^p)).
\]
The $G$-acyclicity hypothesis says precisely that this term vanishes for
$q>0$, while for $q=0$ it is $GF(I^p)$. Thus this spectral sequence has
only one nonzero row and degenerates. Its total cohomology is therefore
\[
	\Hm^n(GF(I^\bullet))=\mathrm{R}^n(GF)(X).
\]

For the second filtration, the defining property of a Cartan--Eilenberg
resolution says that the complex obtained after taking cohomology in the
$F(I^\bullet)$ direction is an injective resolution of
\[
	\Hm^q(F(I^\bullet))=\mathrm{R}^qF(X).
\]
After applying $G$ and taking cohomology, one consequently obtains
\[
	E_2^{p,q}=\mathrm{R}^pG(\mathrm{R}^qF(X)).
\]
Both filtrations compute the same total cohomology, which the first
calculation identifies as $\mathrm{R}^{p+q}(GF)(X)$. Since the double
complex lies in the first quadrant, the filtration in every total degree is
finite and convergence follows. This is the construction in
Theorem~\ref{prop:Grothendieck-ss}; without $G$-acyclicity, the first
spectral sequence need not collapse, and this argument cannot identify the
abutment with the derived functors of the composite.

% stable-exercise-id: o014.mastery.derived-spectral.002.ex06
\subsubsection*{Exercise 6 --- Edge maps and the Grothendieck five-term sequence}
\hypertarget{o014.mastery.derived-spectral.002.ex06}{ }
\label{mastery:derived-spectral-002-ex06}

Use the spectral sequence from Exercise 5.
\begin{enumerate}[(i)]
	\item Identify the two edge morphisms
	\[
		\mathrm{R}^nG(FX)\longrightarrow\mathrm{R}^n(GF)(X),
		\qquad
		\mathrm{R}^n(GF)(X)\longrightarrow G(\mathrm{R}^nF(X)).
	\]
	\item Derive the five-term exact sequence
	\[
	\begin{aligned}
		0\to\mathrm{R}^1G(FX)&\to\mathrm{R}^1(GF)(X)
		\to G(\mathrm{R}^1F(X))\\
		&\xrightarrow{d_2}\mathrm{R}^2G(FX)
		\to\mathrm{R}^2(GF)(X).
	\end{aligned}
	\]
	\item Prove that the first edge morphism is an isomorphism for every
	$n$ if $\mathrm{R}^qF(X)=0$ for $q>0$, and that the second edge
	morphism is an isomorphism for every $n$ if
	$\mathrm{R}^pG(\mathrm{R}^qF(X))=0$ for $p>0$.
\end{enumerate}

% stable-solution-id: o014.mastery.derived-spectral.002.sol06
\paragraph{Solution.}
\hypertarget{o014.mastery.derived-spectral.002.sol06}{ }
Along the bottom edge,
$E_2^{n,0}=\mathrm{R}^nG(FX)$ maps surjectively, page by page, to
$E_\infty^{n,0}$. The latter term is the deepest filtration piece
$\mathrm{F}^n\mathrm{R}^n(GF)(X)$, and its inclusion in the abutment gives
the first edge morphism. Along the left edge, the quotient
\[
	\mathrm{R}^n(GF)(X)/\mathrm{F}^1
	\simeq E_\infty^{0,n}
\]
embeds in $E_2^{0,n}=G(\mathrm{R}^nF(X))$; their composite gives the
second edge morphism. These directions agree with
\eqref{eqn:edge-ss-coh}.

The computation in the working method above applies to every first-quadrant
spectral sequence and gives
\[
	0\to E_2^{1,0}\to H^1\to E_2^{0,1}
	\xrightarrow{d_2}E_2^{2,0}\to H^2.
\]
Substituting
\[
	E_2^{p,q}=\mathrm{R}^pG(\mathrm{R}^qF(X)),
	\qquad H^n=\mathrm{R}^n(GF)(X)
\]
gives exactly the stated five-term sequence.

If $\mathrm{R}^qF(X)=0$ for $q>0$, only the row $q=0$ is nonzero. There
are no differentials or cross-row extension problems, so
\[
	E_2^{n,0}=E_\infty^{n,0}=\mathrm{R}^n(GF)(X),
\]
and the first edge morphism is an isomorphism. If all terms with $p>0$
vanish, only the column $p=0$ remains; similarly,
\[
	\mathrm{R}^n(GF)(X)=E_\infty^{0,n}=E_2^{0,n}
	=G(\mathrm{R}^nF(X)),
\]
and the second edge morphism is an isomorphism.

% stable-exercise-id: o014.mastery.derived-spectral.002.ex07
\subsubsection*{Exercise 7 --- Cohomology of a cyclic group in concrete terms}
\hypertarget{o014.mastery.derived-spectral.002.ex07}{ }
\label{mastery:derived-spectral-002-ex07}

Let $C_m=\lrangle{\sigma}$ be the cyclic group of order $m$, and equip
$A=\Z/r\Z$ with the trivial $C_m$-action. Write $d=\gcd(m,r)$. Use the
periodic resolution
\[
	\cdots\xrightarrow{\nu}\Z[C_m]
	\xrightarrow{\sigma-1}\Z[C_m]
	\xrightarrow{\nu}\Z[C_m]
	\xrightarrow{\sigma-1}\Z[C_m]\to\Z\to0,
	\qquad \nu=1+\sigma+\cdots+\sigma^{m-1},
\]
to compute $\Hm^n(C_m,A)$ for all $n\geq0$. As a check, also determine the
answer for the trivial coefficients $A=\Z$.

% stable-solution-id: o014.mastery.derived-spectral.002.sol07
\paragraph{Solution.}
\hypertarget{o014.mastery.derived-spectral.002.sol07}{ }
Every $\Hom_{\Z[C_m]}(\Z[C_m],A)$ is canonically isomorphic to $A$.
Because the action is trivial, $\sigma-1$ acts as zero, while $\nu$ acts as
multiplication by $m$. Applying $\Hom_{\Z[C_m]}(-,A)$ therefore gives the
complex
\[
	0\to\underbracket{A}_{0}\xrightarrow{0}
	\underbracket{A}_{1}\xrightarrow{\times m}A
	\xrightarrow{0}A\xrightarrow{\times m}A\to\cdots.
\]
Consequently,
\[
	\Hm^n(C_m,A)\simeq
	\begin{cases}
		A,&n=0,\\
		A[m]:=\Ker(\times m:A\to A),&n>0\text{ odd},\\
		A/mA,&n>0\text{ even}.
	\end{cases}
\]
For $A=\Z/r\Z$, the subgroup $A[m]$ is generated by the class of $r/d$
and has order $d$, so $A[m]\simeq\Z/d\Z$. The subgroup $mA$ has order
$r/d$, so the quotient $A/mA$ is also cyclic of order $d$. Thus
\[
	\Hm^n(C_m,\Z/r\Z)\simeq
	\begin{cases}
		\Z/r\Z,&n=0,\\
		\Z/d\Z,&n>0.
	\end{cases}
\]
For $A=\Z$, the kernel of multiplication by $m$ is zero and its cokernel is
$\Z/m\Z$. Hence
\[
	\Hm^n(C_m,\Z)\simeq
	\begin{cases}
		\Z,&n=0,\\
		0,&n>0\text{ odd},\\
		\Z/m\Z,&n>0\text{ even}.
	\end{cases}
\]
This is the direct specialization of
Proposition~\ref{prop:group-cohomology-cyclic}.

% stable-exercise-id: o014.mastery.derived-spectral.002.ex08
\subsubsection*{Exercise 8 --- Checking a multiplicative filtration}
\hypertarget{o014.mastery.derived-spectral.002.ex08}{ }
\label{mastery:derived-spectral-002-ex08}

Let $(A^\bullet,d,\mu)$ be a cohomological dg algebra over a commutative
ring $\Bbbk$, equipped with a decreasing filtration that is exhaustive and
bounded in every degree,
\[
	d(\mathrm{F}^pA)\subset\mathrm{F}^pA,
	\qquad
	\mathrm{F}^pA\cdot\mathrm{F}^{p'}A
	\subset\mathrm{F}^{p+p'}A.
\]
\begin{enumerate}[(i)]
	\item Prove that $E_0=\gr_{\mathrm F}A$ has a graded multiplication and
	that $d_0$ satisfies the Leibniz rule.
	\item If $x$ and $y$ represent classes in
	$E_r^{p,q}$ and $E_r^{p',q'}$, prove that $xy$ represents a class in
	$E_r^{p+p',q+q'}$, that its product is independent of the chosen
	representatives, and that
	\[
		d_r(xy)=d_r(x)y+(-1)^{p+q}x\,d_r(y).
	\]
	\item For the induced filtration
	$\mathrm{F}^p\Hm^n(A)=\Image[\Hm^n(\mathrm{F}^pA)\to\Hm^n(A)]$,
	prove that
	\[
		\mathrm{F}^p\Hm^n(A)\cdot\mathrm{F}^{p'}\Hm^{n'}(A)
		\subset\mathrm{F}^{p+p'}\Hm^{n+n'}(A)
	\]
	and that the convergence isomorphism
	$E_\infty\simeq\gr_{\mathrm F}\Hm(A)$ preserves multiplication.
\end{enumerate}

% stable-solution-id: o014.mastery.derived-spectral.002.sol08
\paragraph{Solution.}
\hypertarget{o014.mastery.derived-spectral.002.sol08}{ }
For homogeneous classes
$[x]\in\gr^p_{\mathrm F}A^n$ and
$[y]\in\gr^{p'}_{\mathrm F}A^{n'}$, define
\[
	[x][y]:=[xy]\in\gr^{p+p'}_{\mathrm F}A^{n+n'}.
\]
If $x$ or $y$ is changed by an element at the next filtration level,
multiplicative compatibility places the resulting change in $xy$ inside
$\mathrm{F}^{p+p'+1}$, so the multiplication is well defined. Since $d$
preserves the filtration, it induces $d_0$, and the Leibniz rule in $A$
gives
\[
	d_0([x][y])=d_0([x])\,[y]+(-1)^n[x]\,d_0([y]).
\]

On the $E_r$ page, representatives can be chosen so that
\[
	x\in\mathrm{F}^pA^{p+q},\quad
	dx\in\mathrm{F}^{p+r}A^{p+q+1},
\]
and similarly
$y\in\mathrm{F}^{p'}A^{p'+q'}$ with
$dy\in\mathrm{F}^{p'+r}A^{p'+q'+1}$. Filtration compatibility and the
Leibniz rule give
\[
	xy\in\mathrm{F}^{p+p'}A^{p+q+p'+q'},
\]
and
\[
	d(xy)=dx\,y+(-1)^{p+q}x\,dy
	\in\mathrm{F}^{p+p'+r}A^{p+q+p'+q'+1}.
\]
Thus $xy$ indeed determines a class in $E_r^{p+p',q+q'}$, and the Leibniz
formula for $d_r$ follows immediately.

To check that the result is independent of the representatives, it is
enough to consider changes by $r$-boundaries and by elements in a deeper
filtration. If, for example, $x$ is replaced by $x+du$, with
$u\in\mathrm{F}^{p-r+1}A^{p+q-1}$, then
\[
	(du)y=d(uy)-(-1)^{p+q-1}u\,dy.
\]
The first term is a boundary, while the second lies in
\[
	\mathrm{F}^{p-r+1+p'+r}A
	=\mathrm{F}^{p+p'+1}A,
\]
so it does not change the class in question. A change of the representative
$y$ is handled symmetrically; changes in a deeper filtration are even more
immediate. Thus multiplication is well defined on every page and passes to
$E_{r+1}\simeq\Hm(E_r,d_r)$.

Finally, if $[x]\in\mathrm{F}^p\Hm^n(A)$ and
$[y]\in\mathrm{F}^{p'}\Hm^{n'}(A)$, choose closed representatives
$x\in\mathrm{F}^pA^n$ and $y\in\mathrm{F}^{p'}A^{n'}$. Their product
$xy$ is closed and lies in $\mathrm{F}^{p+p'}A^{n+n'}$, which proves the
required filtration inclusion. The isomorphism
$E_\infty^{p,q}\simeq\gr^p_{\mathrm F}\Hm^{p+q}(A)$ is built from the
same filtered representatives on both sides. The product of two
representatives maps to the class of their product, so this isomorphism
preserves multiplication. This check is the argument represented in
Proposition~\ref{prop:dga-mult-ss}.
